% ========================================================
% BÀI TẬP TOÁN 5 - KẾT NỐI TRI THỨC VỚI CUỘC SỐNG
% Bài 68: Ôn tập số tự nhiên, phân số, số thập phân
% Định dạng: BTTeX Part III - Trắc nghiệm điền số ngắn
% ========================================================

% Câu 1
\begin{ex}
  Trong số $345\,678$, chữ số $4$ thuộc hàng nào? (Điền số: Hàng đơn vị = 1, Chục = 10, Trăm = 100, Nghìn = 1000, Chục nghìn = 10000, Trăm nghìn = 100000)
  \shortans{10000}
  \loigiai{
    Chữ số 4 thuộc hàng chục nghìn (10 000).
  }
\end{ex}

% Câu 2
\begin{ex}
  Giá trị của chữ số $5$ trong số $345\,678$ là bao nhiêu nghìn?
  \shortans{5}
  \loigiai{
    Chữ số 5 có giá trị 5 000, tức 5 nghìn.
  }
\end{ex}

% Câu 3
\begin{ex}
  Rút gọn phân số $\frac{24}{36}$ về phân số tối giản. Mẫu số của phân số tối giản là bao nhiêu?
  \shortans{3}
  \loigiai{
    $\frac{24}{36} = \frac{2}{3}$. Mẫu số là 3.
  }
\end{ex}

% Câu 4
\begin{ex}
  Ở câu trên, tử số của phân số tối giản là bao nhiêu?
  \shortans{2}
  \loigiai{
    Tử số là 2.
  }
\end{ex}

% Câu 5
\begin{ex}
  Quy đồng mẫu số hai phân số $\frac{3}{4}$ và $\frac{5}{6}$ với mẫu số chung nhỏ nhất là bao nhiêu?
  \shortans{12}
  \loigiai{
    MSCNN(4, 6) = 12.
  }
\end{ex}

% Câu 6
\begin{ex}
  Trong số thập phân $85{,}462$, chữ số $6$ thuộc hàng phần mấy? (Điền số: phần mười = 10, phần trăm = 100, phần nghìn = 1000)
  \shortans{100}
  \loigiai{
    Chữ số 6 thuộc hàng phần trăm (100).
  }
\end{ex}

% Câu 7
\begin{ex}
  Giá trị của chữ số $4$ trong số thập phân $85{,}462$ khi nhân với 10 là bao nhiêu?
  \shortans{4}
  \loigiai{
    Chữ số 4 có giá trị là $0{,}4$. Nhân với 10 được 4.
  }
\end{ex}

% Câu 8
\begin{ex}
  Viết hỗn số $3\frac{2}{5}$ dưới dạng số thập phân. Số thập phân đó nhân với 10 được bao nhiêu?
  \shortans{34}
  \loigiai{
    $3\frac{2}{5} = 3{,}4$. Nhân 10 được 34.
  }
\end{ex}

% Câu 9
\begin{ex}
  Viết phân số $\frac{3}{4}$ dưới dạng số thập phân. Số thập phân đó nhân với 100 được bao nhiêu?
  \shortans{75}
  \loigiai{
    $\frac{3}{4} = 0{,}75$. Nhân 100 được 75.
  }
\end{ex}

% Câu 10
\begin{ex}
  Làm tròn số thập phân $45{,}68$ đến hàng đơn vị được số nào?
  \shortans{46}
  \loigiai{
    $45{,}68 \approx 46$.
  }
\end{ex}

% Câu 11
\begin{ex}
  Làm tròn số thập phân $12{,}345$ đến hàng phần mười. Kết quả nhân với 10 được bao nhiêu?
  \shortans{123}
  \loigiai{
    $12{,}345 \approx 12{,}3$. Nhân 10 được 123.
  }
\end{ex}

% Câu 12
\begin{ex}
  Tìm chữ số $x$ để số $\overline{45x}$ chia hết cho cả 2 và 5.
  \shortans{0}
  \loigiai{
    Số chia hết cho cả 2 và 5 tận cùng là 0. Vậy $x = 0$.
  }
\end{ex}

% Câu 13
\begin{ex}
  Tìm chữ số $x$ để số $\overline{23x}$ chia hết cho 9.
  \shortans{4}
  \loigiai{
    $2 + 3 + x = 5 + x$ chia hết cho 9 nên $x = 4$.
  }
\end{ex}

% Câu 14
\begin{ex}
  Có bao nhiêu số tự nhiên chẵn có 2 chữ số?
  \shortans{45}
  \loigiai{
    Các số từ 10 đến 98: $(98 - 10) : 2 + 1 = 45$ số.
  }
\end{ex}

% Câu 15
\begin{ex}
  Có bao nhiêu phân số có mẫu số là 10 lớn hơn 0 và bé hơn 1?
  \shortans{9}
  \loigiai{
    Các phân số từ $\frac{1}{10}$ đến $\frac{9}{10}$: có 9 phân số.
  }
\end{ex}

% Câu 16
\begin{ex}
  So sánh hai phân số $\frac{5}{8}$ và $\frac{7}{8}$. Phân số lớn hơn có tử số là bao nhiêu?
  \shortans{7}
  \loigiai{
    $\frac{7}{8} > \frac{5}{8}$. Tử số lớn hơn là 7.
  }
\end{ex}

% Câu 17
\begin{ex}
  Trong các số $4{,}56; 4{,}65; 4{,}605; 4{,}506$, số lớn nhất khi nhân với 100 là bao nhiêu?
  \shortans{465}
  \loigiai{
    Số lớn nhất là $4{,}65$. Nhân 100 được 465.
  }
\end{ex}

% Câu 18
\begin{ex}
  Ở các số trên, số bé nhất khi nhân với 1\,000 là bao nhiêu?
  \shortans{4506}
  \loigiai{
    Số bé nhất là $4{,}506$. Nhân 1\,000 được 4506.
  }
\end{ex}

% Câu 19
\begin{ex}
  Viết số thập phân $0{,}08$ dưới dạng phân số tối giản. Mẫu số của phân số đó là bao nhiêu?
  \shortans{25}
  \loigiai{
    $0{,}08 = \frac{8}{100} = \frac{2}{25}$. Mẫu số là 25.
  }
\end{ex}

% Câu 20
\begin{ex}
  Ở câu trên, tử số là bao nhiêu?
  \shortans{2}
  \loigiai{
    Tử số là 2.
  }
\end{ex}

% Câu 21
\begin{ex}
  Một lớp học có 18 học sinh nam và 24 học sinh nữ. Rút gọn tỉ số giữa số nam và nữ về phân số tối giản. Mẫu số là bao nhiêu?
  \shortans{4}
  \loigiai{
    $\frac{18}{24} = \frac{3}{4}$. Mẫu số là 4.
  }
\end{ex}

% Câu 22
\begin{ex}
  Số tự nhiên nhỏ nhất có 4 chữ số khác nhau là bao nhiêu?
  \shortans{1023}
  \loigiai{
    Số đó là 1023.
  }
\end{ex}

% Câu 23
\begin{ex}
  Số tự nhiên lớn nhất có 4 chữ số khác nhau là bao nhiêu?
  \shortans{9876}
  \loigiai{
    Số đó là 9876.
  }
\end{ex}

% Câu 24
\begin{ex}
  Tìm phần nguyên của số thập phân $234{,}567$.
  \shortans{234}
  \loigiai{
    Phần nguyên là 234.
  }
\end{ex}

% Câu 25
\begin{ex}
  Cho 4 chữ số 0, 1, 2, 3. Viết được bao nhiêu số tự nhiên có 3 chữ số khác nhau?
  \shortans{18}
  \loigiai{
    Hàng trăm có 3 cách (1,2,3); hàng chục có 3 cách; hàng đơn vị có 2 cách. $3 \times 3 \times 2 = 18$ số.
  }
\end{ex}
